\chapter{Merz-Wüthrich: the MSEP of the one-year claims development result}

Merz and Wüthrich (2008), \emph{Modelling the claims development result for solvency purposes},
Casualty Actuarial Society E-Forum, Fall 2008, pages 542-568, quantify the uncertainty of the
one-year move of the predicted ultimate. The previous chapter defined the true claims development
result and proved that it has conditional mean zero; this chapter adds its second moment, which is
the conditional mean squared error of predicting it by the budget value $0$, and records the
estimator the paper displays for the observable claims development result.

Index conventions. The paper writes $I$ for the last accident year, $J = I$ for the last
development year, and $\mathcal{D}_I$ for the triangle after $I$ diagonals. Here $I = J = n-1$;
$I - i$ is $n-1-i$, $C_{I-j,j}$ is $C_{n-1-j,\,j}$, $S^I_j$ is $S_j$ computed with $n$ accident
years and $S^{I+1}_j$ is the same column sum computed with $n+1$, that is, with the next diagonal
added. Both are functions of the triangle observed today, because
$S^{I+1}_j = S^I_j + C_{I-j,j}$ and $C_{I-j,j}$ sits on the latest observed diagonal.

The step formalized below is the development-year step $\mathcal{D}_k \to \mathcal{D}_{k+1}$ of
this library's filtration, not the paper's calendar-year step; the caveat is the one recorded for
Theorem~\ref{thm:cdr_closed}.

\section{The exact one-year process variance}

\begin{definition}[One-year process variance]\label{def:cdrProcVar}\lean{VerifiedReserving.cdrProcVar}\leanok
\uses{def:Mack3}
For claims $c$ at development year $k$ and ultimate development year $N$,
\[
  \mathrm{cdrProcVar}(c, f, \sigma^2, N, k) = \Big(\prod_{j=k+1}^{N-1} f_j\Big)^2 \sigma_k^2\, c .
\]
\end{definition}

\begin{theorem}[Merz-Wüthrich 2008, display (2.18)]\label{thm:cdr_msep_exact}\lean{VerifiedReserving.condExp_sq_trueCDR, VerifiedReserving.condVar_trueCDR}\leanok
\uses{def:cdrProcVar, def:trueCDR, thm:cdr_closed, thm:cdr_martingale, def:Mack1, def:Mack3}
Under (M1) and (M3), for $i < n$ and $k < n-1$ and integrable claims,
\[
  E\big[\mathrm{CDR}_i(k)^2 \mid \mathcal{D}_k\big]
   = \Big(\prod_{j=k+1}^{n-2} f_j\Big)^2 \sigma_k^2\, C_{i,k}
\]
almost surely. Since the true CDR has conditional mean zero, this is at once its conditional
variance and the conditional mean squared error of predicting it by $0$, which is the budget value
the paper enters in the income statement for prior accident years.
\end{theorem}
\begin{proof}\leanok
By Theorem~\ref{thm:cdr_closed} the true CDR is $\big(\prod_{j>k} f_j\big)$ times the one-step
residual $\varepsilon_{i,k} = C_{i,k+1} - f_k C_{i,k}$, up to sign. The constant comes out of the
conditional expectation squared, and $E[\varepsilon_{i,k}^2 \mid \mathcal{D}_k] = \sigma_k^2 C_{i,k}$
is assumption (M3). The conditional-variance form subtracts
Theorem~\ref{thm:cdr_martingale} squared, which is $0$.
\end{proof}

\section{The estimator of Merz and Wüthrich}

Definitions only. The derivation in the paper's appendix replaces the estimated development factors
by resampled ones and drops terms of higher order in the residuals, so none of the objects below is
upgraded to an identity.

\begin{lemma}[The column sum one year later]\label{lem:S_succ}\lean{VerifiedReserving.S_succ_eq}\leanok
\uses{def:S}
For $j < n$, $S^{I+1}_j = S^I_j + C_{I-j,j}$.
\end{lemma}
\begin{proof}\leanok The index set gains exactly the row $I-j$. \end{proof}

\begin{definition}[Displays (3.4)-(3.7)]\label{def:mw}\lean{VerifiedReserving.mwPsi, VerifiedReserving.mwPhi, VerifiedReserving.mwGamma, VerifiedReserving.mwDelta}\leanok
\uses{def:sigma2, def:fhat, def:S, lem:S_succ}
With $d = I - i$,
\begin{align*}
 \hat\Psi^I_i &= \frac{\hat\sigma^2_{d}/(\hat f^I_{d})^2}{C_{i,d}}, &
 \hat\Phi^I_{i,J} &= \sum_{j=d+1}^{J-1}\Big(\frac{C_{I-j,j}}{S^{I+1}_j}\Big)^2
   \frac{\hat\sigma^2_j/(\hat f^I_j)^2}{C_{I-j,j}}, \\
 \hat\Gamma^I_{i,J} &= \hat\Phi^I_{i,J} + \hat\Psi^I_i, &
 \hat\Delta^I_{i,J} &= \frac{\hat\sigma^2_{d}/(\hat f^I_{d})^2}{S^I_{d}}
   + \sum_{j=d+1}^{J-1}\Big(\frac{C_{I-j,j}}{S^{I+1}_j}\Big)^2
     \frac{\hat\sigma^2_j/(\hat f^I_j)^2}{S^I_j}.
\end{align*}
The later development steps enter the one-year figure only through the re-estimation of their
development factors on the next diagonal, which is why each is scaled by
$(C_{I-j,j}/S^{I+1}_j)^2 \le 1$.
\end{definition}

\begin{definition}[Displays (3.8)-(3.10)]\label{def:mwMsep}\lean{VerifiedReserving.mwVarCDR, VerifiedReserving.mwProcessCDR, VerifiedReserving.mwParamCDR, VerifiedReserving.mwMsepCDR, VerifiedReserving.mwMsepObsCDR}\leanok
\uses{def:mw, def:reserve}
$\widehat{\mathrm{Var}}(\mathrm{CDR}_i(I+1) \mid \mathcal{D}_I) = (\hat C^I_{i,J})^2 \hat\Psi^I_i$;
the prospective estimator
$\widehat{\mathrm{msep}}(0) = (\hat C^I_{i,J})^2(\hat\Gamma^I_{i,J} + \hat\Delta^I_{i,J})$ split
into its process part $(\hat C^I_{i,J})^2\hat\Gamma^I_{i,J}$ and its parameter part
$(\hat C^I_{i,J})^2\hat\Delta^I_{i,J}$; and the retrospective estimator
$\widehat{\mathrm{msep}}(\widehat{\mathrm{CDR}}_i(I+1)) = (\hat C^I_{i,J})^2(\hat\Phi^I_{i,J} + \hat\Delta^I_{i,J})$.
\end{definition}

\begin{lemma}[Display (3.11)]\label{lem:mw_311}\lean{VerifiedReserving.mwMsepCDR_eq_mwMsepObsCDR_add_mwVarCDR, VerifiedReserving.mwMsepObsCDR_le_mwMsepCDR}\leanok
\uses{def:mwMsep}
$\widehat{\mathrm{msep}}(0) = \widehat{\mathrm{msep}}(\widehat{\mathrm{CDR}}_i(I+1))
 + \widehat{\mathrm{Var}}(\mathrm{CDR}_i(I+1)\mid\mathcal{D}_I)$, and the second is at most the
first when $\hat\Psi^I_i \ge 0$.
\end{lemma}
\begin{proof}\leanok $\hat\Gamma = \hat\Phi + \hat\Psi$. \end{proof}

\begin{theorem}[Display (3.17)]\label{thm:mw_display}\lean{VerifiedReserving.mw_term_split, VerifiedReserving.mwMsepCDR_eq_display}\leanok
\uses{def:mwMsep, lem:S_succ}
If the column sums $S^I_j$ and $S^{I+1}_j$ along the row do not vanish, then
\[
 \widehat{\mathrm{msep}}(0) = (\hat C^I_{i,J})^2\Big(
  \frac{\hat\sigma^2_{d}/(\hat f^I_{d})^2}{C_{i,d}}
  + \frac{\hat\sigma^2_{d}/(\hat f^I_{d})^2}{S^I_{d}}
  + \sum_{j=d+1}^{J-1}\frac{C_{I-j,j}}{S^{I+1}_j}\,
     \frac{\hat\sigma^2_j/(\hat f^I_j)^2}{S^I_j}\Big).
\]
Read next to Mack's Theorem 3, the one-year figure keeps only the first term of the process
variance, keeps the estimation error of the next diagonal in full, and scales every remaining
run-off cell by $C_{I-j,j}/S^{I+1}_j \le 1$.
\end{theorem}
\begin{proof}\leanok
The two sums combine term by term: with $s = S^I_j$, $c = C_{I-j,j}$ and $s + c = S^{I+1}_j$ by
Lemma~\ref{lem:S_succ}, $(c/(s+c))^2 a/c + (c/(s+c))^2 a/s = (c/(s+c))\,a/s$.
\end{proof}

\section{The bridge to the exact statement}

\begin{theorem}[The first development step is the plug-in]\label{thm:mw_plugin}\lean{VerifiedReserving.mwVarCDR_eq_plugin, VerifiedReserving.mwProcessCDR_eq_plugin}\leanok
\uses{def:mwMsep, def:cdrProcVar, thm:cdr_msep_exact, def:Chat}
If the row is not fully developed and $\hat f^I_{d} \neq 0$, $C_{i,d} \neq 0$, then
$(\hat C^I_{i,J})^2\hat\Psi^I_i$ is exactly
$\mathrm{cdrProcVar}(C_{i,d}, \hat f, \hat\sigma^2, n-1, d)$: display (3.8) is the plug-in of
Theorem~\ref{thm:cdr_msep_exact} for the step the next calendar year observes. The process part of
display (3.9) is that plug-in plus the correction $(\hat C^I_{i,J})^2\hat\Phi^I_{i,J}$ for the
later steps.
\end{theorem}
\begin{proof}\leanok
Split the ultimate at its bottom factor, $\hat C^I_{i,J} = C_{i,d}\,\hat f^I_d \prod_{j>d}\hat f^I_j$;
the squares of $C_{i,d}$ and $\hat f^I_d$ cancel against the denominators of $\hat\Psi^I_i$.
\end{proof}

\begin{theorem}[The one-year process variances telescope]\label{thm:cdr_telescope}\lean{VerifiedReserving.sum_range_cdrProcVar_eq_procVar, VerifiedReserving.sum_procVar_step_eq_procVar, VerifiedReserving.condExp_sum_cdrProcVar_eq_procVar}\leanok
\uses{def:cdrProcVar, def:procVar, lem:C_at}
Summing the exact one-year process variance of Theorem~\ref{thm:cdr_msep_exact} over every
remaining development step of the row, with $C_{i,k}$ replaced by its (M1) prediction
$C_{i,d}\prod_{l=d}^{k-1} f_l$, gives Mack's multi-year process variance $V_m$ exactly, with no
remainder. The stochastic form performs the replacement by a conditional expectation given
$\mathcal{D}_d$ instead of by hand.
\end{theorem}
\begin{proof}\leanok
Term by term this is the closed form of $V_m$ (Definition~\ref{def:procVar}): its $j$-th summand is
the one-year process variance of step $d+j$. For the stochastic form, each summand is a constant
times $C_{i,d+j}$, and Lemma~\ref{lem:C_at} evaluates its conditional expectation.
\end{proof}

\section{Aggregation over accident years}

Accident year $0$ is fully developed, so the sums run over $i = 1, \dots, n-1$ and the cross terms
over the pairs $k > i > 0$, the older year $i$ carrying the shared development factors.
Definitions only.

\begin{definition}[Displays (3.12)-(3.15)]\label{def:mw_total}\lean{VerifiedReserving.mwLambda, VerifiedReserving.mwXi, VerifiedReserving.mwCrossObs, VerifiedReserving.mwCross, VerifiedReserving.mwMsepTotalObsCDR, VerifiedReserving.mwMsepTotalCDR}\leanok
\uses{def:mwMsep, def:total}
\[
 \hat\Lambda^I_{i,J} = \frac{C_{i,d}}{S^{I+1}_{d}}\,\frac{\hat\sigma^2_{d}/(\hat f^I_{d})^2}{S^I_{d}}
  + \sum_{j=d+1}^{J-1}\Big(\frac{C_{I-j,j}}{S^{I+1}_j}\Big)^2\frac{\hat\sigma^2_j/(\hat f^I_j)^2}{S^I_j},
 \qquad
 \hat\Xi^I_{i,J} = \hat\Phi^I_{i,J} + \frac{\hat\sigma^2_{d}/(\hat f^I_{d})^2}{S^{I+1}_{d}},
\]
and the two aggregates
$\sum_i \widehat{\mathrm{msep}}_i + 2\sum_{k>i>0}\hat C^I_{i,J}\hat C^I_{k,J}(\cdot_i + \hat\Lambda^I_{i,J})$
with $\cdot_i = \hat\Phi^I_{i,J}$ for the retrospective view (3.12) and $\cdot_i = \hat\Xi^I_{i,J}$
for the prospective view (3.15).
\end{definition}

\begin{theorem}[Display (3.16)]\label{thm:mw_total}\lean{VerifiedReserving.mwMsepTotalCDR_sub_mwMsepTotalObsCDR}\leanok
\uses{def:mw_total, lem:mw_311}
The aggregated prospective estimator exceeds the aggregated retrospective one by
$\sum_i \widehat{\mathrm{Var}}(\mathrm{CDR}_i(I+1)\mid\mathcal{D}_I)
 + 2\sum_{k>i>0}\hat C^I_{i,J}\hat C^I_{k,J}(\hat\Xi^I_{i,J} - \hat\Phi^I_{i,J})$: the same
decoupling for aggregated accident years as for a single one.
\end{theorem}
\begin{proof}\leanok Termwise, Lemma~\ref{lem:mw_311} and ring algebra. \end{proof}
